\documentclass[11pt]{article}

\usepackage{amsmath,amssymb,amsthm,bm,mathtools}
\usepackage{geometry}
\geometry{margin=1in}
\usepackage{setspace}
\usepackage{enumitem}
\usepackage{booktabs}
\usepackage{hyperref}
\usepackage{algorithm}
\usepackage{algpseudocode}
\usepackage{float}

\onehalfspacing

\newcommand{\R}{\mathbb R}
\newcommand{\ind}{\mathbb I}
\newcommand{\norm}[1]{\left\lVert #1 \right\rVert}
\newcommand{\trans}{^{\mathsf T}}
\DeclareMathOperator{\logit}{logit}
\DeclareMathOperator{\Cat}{Categorical}
\DeclareMathOperator{\Dir}{Dirichlet}
\DeclareMathOperator{\IG}{IG}
\DeclareMathOperator{\IW}{IW}
\DeclareMathOperator{\tr}{tr}

\title{Joint Multilayered LSIRM with Latent-Position Gaussian Mixture Clustering\\
NIW Predictive Allocation and Jain--Neal Split--Merge Moves}
\date{Model Description}

\begin{document}
\maketitle

\section{Overview}

This document defines a joint multilayered LSIRM in which item clustering is
performed directly on item latent positions.  The previous PPCA measurement
block is removed.  In the modified model, each item latent position is assigned
a finite Gaussian mixture prior,
\[
    b_j^{(\ell(j))}\mid c_j=l,\mu_l,\Sigma_l
    \sim N_d(\mu_l,\Sigma_l),
    \qquad d=2,
\]
where $c_j$ is the item cluster label.  The mixture prior is imposed only on
the item latent positions.  It is not imposed directly on respondent ability or
item difficulty parameters.  Therefore, clustering regularizes the latent item
geometry while difficulty and ability remain governed by their own non-clustered
priors.

The model is fully joint: item cluster labels affect the posterior distribution
of item positions through the mixture prior, and item positions affect the
network/response likelihood through the LSIRM distance term.  Consequently,
cluster labels can indirectly affect respondent and item effects through the
posterior dependence induced by the likelihood, but no cluster-specific prior is
placed directly on those scalar effects.

\section{Indexing, dimensions, and observed data}

\begin{center}
\begin{tabular}{lll}
\toprule
Symbol & Meaning & Dimension/range \\
\midrule
$\ell$ & layer index & $\ell=1,\dots,L$ \\
$i$ & respondent/person index & $i=1,\dots,n_\ell$ \\
$j$ & local item index within layer $\ell$ & $j=1,\dots,P_\ell$ \\
$q$ & global item index across all layers & $q=1,\dots,P$, $P=\sum_{\ell=1}^L P_\ell$ \\
$\ell(q)$ & layer containing global item $q$ & $\{1,\dots,L\}$ \\
$j(q)$ & local item index of global item $q$ & $\{1,
\dots,P_{\ell(q)}\}$ \\
$d$ & latent space dimension & $d=2$ in the current implementation \\
$K^\star$ & truncation level / maximum number of mixture components & positive integer \\
$K_+$ & number of occupied mixture components & $K_+=|\{l:n_l>0\}|$ \\
\bottomrule
\end{tabular}
\end{center}

The observed data are denoted by
\[
    Y^{(\ell)}=\{y_{ij}^{(\ell)}:(i,j)\in\mathcal O_\ell\},
\]
where $\mathcal O_\ell$ is the set of observed respondent--item pairs in layer
$\ell$.  If there is no missingness, then
$\mathcal O_\ell=\{1,\dots,n_\ell\}\times\{1,\dots,P_\ell\}$.

\section{LSIRM likelihood}

For respondent $i$ and item $j$ in layer $\ell$, define
\begin{align}
    a_i^{(\ell)} &\in \R^d,
    &\text{respondent latent position,} \\
    b_j^{(\ell)} &\in \R^d,
    &\text{item latent position,} \\
    \theta_i^{(\ell)} &\in \R,
    &\text{respondent ability or activity effect,} \\
    \beta_j^{(\ell)} &\in \R,
    &\text{item intercept, easiness, or difficulty-related effect.}
\end{align}

The layer-specific LSIRM linear predictor is
\begin{equation}
    \eta_{ij}^{(\ell)}
    =
    \theta_i^{(\ell)}
    +\beta_j^{(\ell)}
    -\norm{a_i^{(\ell)}-b_j^{(\ell)}}_2.
    \label{eq:lsirm-linear-predictor}
\end{equation}

For a binary response layer, the likelihood is
\begin{align}
    y_{ij}^{(\ell)}\mid \theta_i^{(\ell)},\beta_j^{(\ell)},a_i^{(\ell)},b_j^{(\ell)}
    &\sim \operatorname{Bernoulli}(p_{ij}^{(\ell)}), \\
    \logit\{p_{ij}^{(\ell)}\}
    &=\eta_{ij}^{(\ell)}.
    \label{eq:bernoulli-likelihood}
\end{align}

More generally, if different layers use different observation models or link
functions, write
\begin{equation}
    y_{ij}^{(\ell)}\mid \eta_{ij}^{(\ell)},\xi_\ell
    \sim f_\ell\{y_{ij}^{(\ell)}\mid \eta_{ij}^{(\ell)},\xi_\ell\},
    \qquad (i,j)\in\mathcal O_\ell,
    \label{eq:generic-layer-likelihood}
\end{equation}
where $\xi_\ell$ denotes additional layer-specific nuisance parameters, if any.
The full LSIRM log-likelihood is
\begin{equation}
    \ell_{\mathrm{LSIRM}}(\Theta_{\mathrm{LSIRM}})
    =
    \sum_{\ell=1}^L
    \sum_{(i,j)\in\mathcal O_\ell}
    \log f_\ell\{y_{ij}^{(\ell)}\mid \eta_{ij}^{(\ell)},\xi_\ell\}.
    \label{eq:lsirm-loglik}
\end{equation}
For the Bernoulli-logit case,
\begin{equation}
    \log f_\ell(y\mid \eta)
    =
    y\eta-\log(1+e^\eta).
\end{equation}

\section{Latent-position Gaussian mixture prior for items}

\subsection{Global item position notation}

Define the global item latent position
\begin{equation}
    z_q
    =
    b_{j(q)}^{(\ell(q))}
    \in\R^d,
    \qquad q=1,\dots,P.
    \label{eq:global-item-position}
\end{equation}
The clustering model is placed directly on $z_q$.

\subsection{Finite Gaussian mixture prior}

For $q=1,\dots,P$ and $l=1,\dots,K^\star$,
\begin{align}
    z_q\mid c_q=l,\mu_l,\Sigma_l
    &\sim N_d(\mu_l,\Sigma_l),
    \label{eq:z-mixture}\\
    c_q\mid \rho
    &\sim \Cat(\rho_1,\dots,\rho_{K^\star}),\\
    \rho
    &\sim \Dir(e_0,\dots,e_0).
    \label{eq:rho-prior}
\end{align}

The mixture component parameters follow a Normal--Inverse-Wishart prior:
\begin{align}
    \Sigma_l
    &\sim \IW(\nu_0,S_0),
    \label{eq:sigma-niw-prior}\\
    \mu_l\mid \Sigma_l
    &\sim N_d\left(m_0,\frac{1}{\kappa_0}\Sigma_l\right),
    \qquad l=1,\dots,K^\star.
    \label{eq:mu-niw-prior}
\end{align}
Here $m_0\in\R^d$, $\kappa_0>0$, $\nu_0>d-1$, and
$S_0\in\R^{d\times d}$ is positive definite.

\paragraph{Important modeling restriction.}
The cluster label $c_q$ is allowed to affect $z_q=b_{j(q)}^{(\ell(q))}$ only.
There is no cluster-specific prior of the form
$\theta_i^{(\ell)}\mid c_q$, $\beta_j^{(\ell)}\mid c_q$, or
$\xi_\ell\mid c_q$.  Thus, clusters represent latent item geometry rather than
marginal difficulty, ability, or layer-specific nuisance effects.

\section{Priors for non-clustering LSIRM parameters}

For each layer $\ell$,
\begin{align}
    a_i^{(\ell)}\mid \sigma_{a,\ell}^2
    &\sim N_d(0,\sigma_{a,\ell}^2 I_d),
    && i=1,\dots,n_\ell,
    \label{eq:a-prior}\\
    \theta_i^{(\ell)}\mid \sigma_{\theta,\ell}^2
    &\sim N(0,\sigma_{\theta,\ell}^2),
    && i=1,\dots,n_\ell,
    \label{eq:theta-prior}\\
    \beta_j^{(\ell)}\mid \sigma_{\beta,\ell}^2
    &\sim N(0,\sigma_{\beta,\ell}^2),
    && j=1,\dots,P_\ell.
    \label{eq:beta-prior}
\end{align}
The item position $b_j^{(\ell)}$ does not receive an additional independent
Gaussian prior.  Its prior is the mixture prior in
\eqref{eq:z-mixture}--\eqref{eq:mu-niw-prior}.

Optional variance hyperpriors may be assigned as
\begin{align}
    \sigma_{a,\ell}^2 &\sim \IG(A_a,B_a),\\
    \sigma_{\theta,\ell}^2 &\sim \IG(A_\theta,B_\theta),\\
    \sigma_{\beta,\ell}^2 &\sim \IG(A_\beta,B_\beta).
\end{align}
If the variances are fixed, the corresponding inverse-gamma updates below are
omitted.

\section{Layer comparability, alignment, and scale diagnostics}
\label{sec:layer-comparability}

Because the mixture prior clusters item positions across all layers, the model
assumes that
\[
    b_j^{(\ell)}\in\R^d
\]
is expressed in a comparable coordinate system across layers.  If layer-specific
latent spaces differ greatly in scale, translation, or orientation, the mixture
prior may cluster layer effects rather than phenotype-related item geometry.

At each stored MCMC iteration, compute the layer-specific item-position scale
\begin{equation}
    s_\ell
    =
    \left\{
    \frac{1}{P_\ell d}
    \sum_{j=1}^{P_\ell}
    \norm{b_j^{(\ell)}-\bar b^{(\ell)}}_2^2
    \right\}^{1/2},
    \qquad
    \bar b^{(\ell)}=\frac{1}{P_\ell}\sum_{j=1}^{P_\ell}b_j^{(\ell)}.
    \label{eq:layer-scale-diagnostic}
\end{equation}
Then monitor
\begin{equation}
    R_{\mathrm{scale}}
    =
    \frac{\max_\ell s_\ell}{\min_\ell s_\ell}.
\end{equation}
Large values of $R_{\mathrm{scale}}$ indicate that cluster labels may be partly
recovering layer-scale differences.  In that case, the model should be revised
by imposing stronger shared-scale identification, fixing layer-specific latent
space scales, or applying a prespecified alignment/scaling constraint during
MCMC.  This check is part of the model diagnostic procedure, not a separate
post-hoc clustering baseline.

For translational non-identifiability, impose or monitor the centering
constraint
\begin{equation}
    \sum_{i=1}^{n_\ell}a_i^{(\ell)}=0,
    \qquad
    \sum_{j=1}^{P_\ell}b_j^{(\ell)}=0,
    \qquad \ell=1,\dots,L,
    \label{eq:centering-constraint}
\end{equation}
or use an equivalent deterministic recentering step after each complete latent
position update.  Rotational and reflection invariance may be handled by fixing
anchors or by a deterministic Procrustes identification step when necessary.

\section{Full joint posterior}

Let
\begin{align*}
    \Theta_{\mathrm{LSIRM}}
    &=
    \{a_i^{(\ell)},b_j^{(\ell)},\theta_i^{(\ell)},\beta_j^{(\ell)},
    \xi_\ell,\sigma_{a,\ell}^2,\sigma_{\theta,\ell}^2,
    \sigma_{\beta,\ell}^2:
    \ell=1,\dots,L\},\\
    \Theta_{\mathrm{mix}}
    &=
    \{c_q,\rho,\mu_l,\Sigma_l:q=1,\dots,P,\ l=1,\dots,K^\star\}.
\end{align*}
The full posterior is proportional to
\begin{align}
&p(\Theta_{\mathrm{LSIRM}},\Theta_{\mathrm{mix}}\mid Y)
\notag\\
&\quad \propto
\prod_{\ell=1}^L
\prod_{(i,j)\in\mathcal O_\ell}
    f_\ell\{y_{ij}^{(\ell)}\mid \eta_{ij}^{(\ell)},\xi_\ell\}
\notag\\
&\qquad\times
\prod_{q=1}^P
    N_d(z_q;\mu_{c_q},\Sigma_{c_q})\rho_{c_q}
\notag\\
&\qquad\times
\prod_{l=1}^{K^\star}
    p(\mu_l\mid\Sigma_l)p(\Sigma_l)
    \ p(\rho)
\notag\\
&\qquad\times
\prod_{\ell=1}^L
\left[
\prod_{i=1}^{n_\ell}p(a_i^{(\ell)}\mid\sigma_{a,\ell}^2)
\prod_{i=1}^{n_\ell}p(\theta_i^{(\ell)}\mid\sigma_{\theta,\ell}^2)
\prod_{j=1}^{P_\ell}p(\beta_j^{(\ell)}\mid\sigma_{\beta,\ell}^2)
\right]
\notag\\
&\qquad\times
\prod_{\ell=1}^L
    p(\sigma_{a,\ell}^2)p(\sigma_{\theta,\ell}^2)p(\sigma_{\beta,\ell}^2)
    p(\xi_\ell),
    \label{eq:full-posterior}
\end{align}
where $z_q=b_{j(q)}^{(\ell(q))}$.

\section{NIW posterior quantities}

For a generic set of global item indices $A\subseteq\{1,\dots,P\}$, define
\begin{align}
    n_A &= |A|,\\
    \bar z_A &= \frac{1}{n_A}\sum_{q\in A}z_q,
    \qquad n_A>0,\\
    Q_A &= \sum_{q\in A}(z_q-\bar z_A)(z_q-\bar z_A)\trans,
    \qquad n_A>0.
\end{align}
For $n_A=0$, use prior quantities directly and set $Q_A=0$.

The NIW posterior parameters for the items in $A$ are
\begin{align}
    \kappa_A &= \kappa_0+n_A,\\
    \nu_A &= \nu_0+n_A,\\
    m_A &=
    \begin{cases}
    \dfrac{\kappa_0m_0+n_A\bar z_A}{\kappa_0+n_A}, & n_A>0,\\[1em]
    m_0, & n_A=0,
    \end{cases}\\
    S_A &=
    \begin{cases}
    S_0+Q_A+
    \dfrac{\kappa_0 n_A}{\kappa_0+n_A}
    (\bar z_A-m_0)(\bar z_A-m_0)\trans,
    & n_A>0,\\[1em]
    S_0, & n_A=0.
    \end{cases}
    \label{eq:niw-posterior-params}
\end{align}

For component $l$, let
\[
    \mathcal J_l=\{q:c_q=l\},
    \qquad n_l=|\mathcal J_l|.
\]
Then
\[
    (\kappa_l,\nu_l,m_l,S_l)
    =
    (\kappa_{\mathcal J_l},\nu_{\mathcal J_l},m_{\mathcal J_l},S_{\mathcal J_l}).
\]
For leave-one-out calculations,
\[
    \mathcal J_{l,-q}=\{r:r\neq q,\ c_r=l\}.
\]

\section{NIW posterior predictive allocation update}

Integrating out $(\mu_l,\Sigma_l)$ for label allocation gives the Student-$t$
posterior predictive density
\begin{equation}
    z_q\mid z_{\mathcal J_{l,-q}}
    \sim
    t_{\nu_{l,-q}-d+1}
    \left(
    m_{l,-q},
    \frac{\kappa_{l,-q}+1}
    {\kappa_{l,-q}(\nu_{l,-q}-d+1)}S_{l,-q}
    \right).
    \label{eq:niw-predictive}
\end{equation}
Therefore the collapsed single-site label update is
\begin{equation}
    P(c_q=l\mid c_{-q},z)
    \propto
    (n_{l,-q}+e_0)
    t_{\nu_{l,-q}-d+1}
    \left(
    z_q;
    m_{l,-q},
    \frac{\kappa_{l,-q}+1}
    {\kappa_{l,-q}(\nu_{l,-q}-d+1)}S_{l,-q}
    \right),
    \label{eq:c-update}
\end{equation}
for $l=1,\dots,K^\star$.  This is a Gibbs update, so the acceptance
probability is one after sampling from the normalized probabilities in
\eqref{eq:c-update}.

\section{Posterior draw of mixture component parameters}

After updating the labels, draw component parameters from their NIW posterior:
\begin{align}
    \Sigma_l\mid z,c
    &\sim \IW(\nu_l,S_l),
    \label{eq:sigma-posterior}\\
    \mu_l\mid \Sigma_l,z,c
    &\sim N_d\left(m_l,\frac{1}{\kappa_l}\Sigma_l\right),
    \qquad l=1,\dots,K^\star.
    \label{eq:mu-posterior}
\end{align}
For empty components, $n_l=0$, and this reduces to a draw from the NIW prior.

The mixture weights may be sampled for monitoring or storage:
\begin{equation}
    \rho\mid c
    \sim
    \Dir(e_0+n_1,\dots,e_0+n_{K^\star}).
    \label{eq:rho-posterior}
\end{equation}
The draw of $\rho$ is not needed for the collapsed label update.

\section{Collapsed partition target for split--merge moves}

Integrating out $\rho$ and all component parameters gives the collapsed
partition target
\begin{equation}
    \pi(c\mid z)
    \propto
    \prod_{l=1}^{K^\star}
    \frac{\Gamma(n_l+e_0)}{\Gamma(e_0)}
    m(z_{\mathcal J_l}),
    \label{eq:collapsed-partition-target}
\end{equation}
where
\begin{equation}
    m(z_A)
    =
    \int
    \prod_{q\in A}N_d(z_q;\mu,\Sigma)
    p(\mu,\Sigma)\,d\mu\,d\Sigma
\end{equation}
is the NIW marginal likelihood.

The log NIW marginal likelihood is
\begin{align}
    \log m(z_A)
    &=-\frac{n_A d}{2}\log\pi
    +\frac{d}{2}\{\log\kappa_0-\log\kappa_A\}
    +\frac{\nu_0}{2}\log|S_0|
    -\frac{\nu_A}{2}\log|S_A|
    \notag\\
    &\qquad
    +\log\Gamma_d\left(\frac{\nu_A}{2}\right)
    -\log\Gamma_d\left(\frac{\nu_0}{2}\right),
    \label{eq:niw-marginal}
\end{align}
where
\begin{equation}
    \log\Gamma_d(a)
    =
    \frac{d(d-1)}{4}\log\pi
    +
    \sum_{h=1}^d\log\Gamma\left(a+\frac{1-h}{2}\right).
\end{equation}
Thus,
\begin{equation}
    \log\pi(c\mid z)
    =
    \sum_{l=1}^{K^\star}
    \left[
    \log\Gamma(n_l+e_0)-\log\Gamma(e_0)
    +\log m(z_{\mathcal J_l})
    \right]
    +\mathrm{constant}.
    \label{eq:log-collapsed-partition-target}
\end{equation}

\section{Metropolis--Hastings updates for LSIRM parameters}

All LSIRM parameters with non-conjugate full conditionals can be updated by
random-walk Metropolis--Hastings.  The formulas below assume symmetric
Gaussian random-walk proposals.  If a non-symmetric proposal is used, the
corresponding proposal-density ratio must be added to the log acceptance ratio.

\subsection{Respondent ability or activity effect $\theta_i^{(\ell)}$}

Propose
\[
    \theta_i^{(\ell)\prime}
    =
    \theta_i^{(\ell)}+\varepsilon_\theta,
    \qquad
    \varepsilon_\theta\sim N(0,s_\theta^2).
\]
Let $\eta_{ij}^{(\ell)\prime}$ be the linear predictor after replacing
$\theta_i^{(\ell)}$ by $\theta_i^{(\ell)\prime}$.  The log acceptance ratio is
\begin{align}
    \log R_\theta
    &=
    \sum_{j:(i,j)\in\mathcal O_\ell}
    \left[
    \log f_\ell\{y_{ij}^{(\ell)}\mid \eta_{ij}^{(\ell)\prime},\xi_\ell\}
    -
    \log f_\ell\{y_{ij}^{(\ell)}\mid \eta_{ij}^{(\ell)},\xi_\ell\}
    \right]
    \notag\\
    &\qquad
    -
    \frac{1}{2\sigma_{\theta,\ell}^2}
    \left\{(\theta_i^{(\ell)\prime})^2-(\theta_i^{(\ell)})^2\right\}.
    \label{eq:theta-mh-ratio}
\end{align}
Accept with probability
\begin{equation}
    \alpha_\theta
    =
    \min\{1,\exp(\log R_\theta)\}.
\end{equation}

\subsection{Item effect $\beta_j^{(\ell)}$}

Propose
\[
    \beta_j^{(\ell)\prime}
    =
    \beta_j^{(\ell)}+\varepsilon_\beta,
    \qquad
    \varepsilon_\beta\sim N(0,s_\beta^2).
\]
The log acceptance ratio is
\begin{align}
    \log R_\beta
    &=
    \sum_{i:(i,j)\in\mathcal O_\ell}
    \left[
    \log f_\ell\{y_{ij}^{(\ell)}\mid \eta_{ij}^{(\ell)\prime},\xi_\ell\}
    -
    \log f_\ell\{y_{ij}^{(\ell)}\mid \eta_{ij}^{(\ell)},\xi_\ell\}
    \right]
    \notag\\
    &\qquad
    -
    \frac{1}{2\sigma_{\beta,\ell}^2}
    \left\{(\beta_j^{(\ell)\prime})^2-(\beta_j^{(\ell)})^2\right\}.
    \label{eq:beta-mh-ratio}
\end{align}
Accept with probability
\begin{equation}
    \alpha_\beta
    =
    \min\{1,\exp(\log R_\beta)\}.
\end{equation}

\subsection{Respondent latent position $a_i^{(\ell)}$}

Propose
\[
    a_i^{(\ell)\prime}
    =
    a_i^{(\ell)}+\varepsilon_a,
    \qquad
    \varepsilon_a\sim N_d(0,s_a^2 I_d).
\]
The log acceptance ratio is
\begin{align}
    \log R_a
    &=
    \sum_{j:(i,j)\in\mathcal O_\ell}
    \left[
    \log f_\ell\{y_{ij}^{(\ell)}\mid \eta_{ij}^{(\ell)\prime},\xi_\ell\}
    -
    \log f_\ell\{y_{ij}^{(\ell)}\mid \eta_{ij}^{(\ell)},\xi_\ell\}
    \right]
    \notag\\
    &\qquad
    -
    \frac{1}{2\sigma_{a,\ell}^2}
    \left\{
    \norm{a_i^{(\ell)\prime}}_2^2-
    \norm{a_i^{(\ell)}}_2^2
    \right\}.
    \label{eq:a-mh-ratio}
\end{align}
Accept with probability
\begin{equation}
    \alpha_a
    =
    \min\{1,\exp(\log R_a)\}.
\end{equation}

\subsection{Item latent position $b_j^{(\ell)}$}

Let $q=q(\ell,j)$ be the global item index corresponding to local item $j$ in
layer $\ell$, and let $c_q$ be its current cluster label.  Propose
\[
    b_j^{(\ell)\prime}
    =
    b_j^{(\ell)}+\varepsilon_b,
    \qquad
    \varepsilon_b\sim N_d(0,s_b^2 I_d).
\]
The log acceptance ratio must include both the LSIRM likelihood contribution
and the Gaussian mixture prior contribution:
\begin{align}
    \log R_b
    &=
    \sum_{i:(i,j)\in\mathcal O_\ell}
    \left[
    \log f_\ell\{y_{ij}^{(\ell)}\mid \eta_{ij}^{(\ell)\prime},\xi_\ell\}
    -
    \log f_\ell\{y_{ij}^{(\ell)}\mid \eta_{ij}^{(\ell)},\xi_\ell\}
    \right]
    \notag\\
    &\qquad
    +
    \log N_d\{b_j^{(\ell)\prime};\mu_{c_q},\Sigma_{c_q}\}
    -
    \log N_d\{b_j^{(\ell)};\mu_{c_q},\Sigma_{c_q}\}.
    \label{eq:b-mh-ratio}
\end{align}
Accept with probability
\begin{equation}
    \alpha_b
    =
    \min\{1,\exp(\log R_b)\}.
    \label{eq:b-acceptance-probability}
\end{equation}
This is the main difference from a post-hoc clustering approach.  The cluster
assignment affects the posterior update of $b_j^{(\ell)}$ through the mixture
prior term in \eqref{eq:b-mh-ratio}.

\section{Conjugate updates for variance parameters}

If the layer-specific prior variances are sampled, their full conditionals are
\begin{align}
    \sigma_{a,\ell}^2\mid a^{(\ell)}
    &\sim
    \IG\left(
    A_a+\frac{n_\ell d}{2},
    B_a+\frac{1}{2}\sum_{i=1}^{n_\ell}\norm{a_i^{(\ell)}}_2^2
    \right),
    \label{eq:sigma-a-update}\\
    \sigma_{\theta,\ell}^2\mid \theta^{(\ell)}
    &\sim
    \IG\left(
    A_\theta+\frac{n_\ell}{2},
    B_\theta+\frac{1}{2}\sum_{i=1}^{n_\ell}(\theta_i^{(\ell)})^2
    \right),
    \label{eq:sigma-theta-update}\\
    \sigma_{\beta,\ell}^2\mid \beta^{(\ell)}
    &\sim
    \IG\left(
    A_\beta+\frac{P_\ell}{2},
    B_\beta+\frac{1}{2}\sum_{j=1}^{P_\ell}(\beta_j^{(\ell)})^2
    \right).
    \label{eq:sigma-beta-update}
\end{align}

\section{Jain--Neal restricted Gibbs split--merge moves}

Split--merge moves are Metropolis--Hastings updates on $c$ conditional on the
current item positions $z=(z_1,\dots,z_P)$.  The target is the collapsed
partition density in \eqref{eq:log-collapsed-partition-target}.

\subsection{Anchor selection}

Select two distinct global item indices $u$ and $v$ uniformly from
$\{1,\dots,P\}$.  If $c_u=c_v$, attempt a split move.  If $c_u\neq c_v$,
attempt a merge move.  The anchor selection probability is symmetric and
cancels in the Metropolis--Hastings ratio.

\subsection{Split proposal}

Suppose $c_u=c_v=a$.  Let
\[
    A=\{q:c_q=a\}
\]
be the cluster to be split.  Choose an empty component label
$b\in\{l:n_l=0\}$ uniformly.  If no empty component exists, skip the split.

Initialize two temporary clusters by forcing
\[
    c_u'=a,
    \qquad
    c_v'=b.
\]
For each $q\in A\setminus\{u,v\}$, assign $q$ to either $a$ or $b$ by a
restricted NIW-collapsed Gibbs scan:
\begin{align}
    w_q(a)&\propto
    (n_{a,-q}^{\mathrm{temp}}+e_0)
    p(z_q\mid z_{a,-q}^{\mathrm{temp}}),\\
    w_q(b)&\propto
    (n_{b,-q}^{\mathrm{temp}}+e_0)
    p(z_q\mid z_{b,-q}^{\mathrm{temp}}),
\end{align}
where the predictive density is the NIW Student-$t$ density in
\eqref{eq:niw-predictive}.  Let $p_q^{\mathrm{fwd}}$ be the normalized
probability of the assignment actually selected during the forward restricted
scan.  Then
\begin{equation}
    \log q(c'\mid c)
    =
    -\log E(c)
    +
    \sum_{q\in A\setminus\{u,v\}}\log p_q^{\mathrm{fwd}},
    \label{eq:split-forward-probability}
\end{equation}
where $E(c)=|\{l:n_l=0\}|$ is the number of empty components before the split.
The reverse move is deterministic merge, so
\[
    \log q(c\mid c')=0.
\]
The split acceptance probability is
\begin{equation}
    \alpha_{\mathrm{split}}
    =
    \min\left\{1,
    \exp\left[
    \log\pi(c'\mid z)-\log\pi(c\mid z)
    -\log q(c'\mid c)
    \right]
    \right\}.
    \label{eq:split-acceptance}
\end{equation}

\subsection{Merge proposal}

Suppose $c_u=a$, $c_v=b$, and $a\neq b$.  Let
\[
    A=\{q:c_q=a\},
    \qquad
    B=\{q:c_q=b\}.
\]
The proposal merges $B$ into $A$:
\[
    c_q'=a,
    \qquad q\in A\cup B.
\]
Thus
\[
    \log q(c'\mid c)=0.
\]
The reverse proposal is a split of $A\cup B$ back into $A$ and $B$ using the
same anchors $u$ and $v$.  Let $p_q^{\mathrm{rev}}$ be the normalized
probability of assigning $q$ to its original side during this reverse
restricted Gibbs scan.  Then
\begin{equation}
    \log q(c\mid c')
    =
    -\log E(c')
    +
    \sum_{q\in(A\cup B)\setminus\{u,v\}}\log p_q^{\mathrm{rev}},
    \label{eq:merge-reverse-probability}
\end{equation}
where $E(c')$ is the number of empty components after the merge.  The merge
acceptance probability is
\begin{equation}
    \alpha_{\mathrm{merge}}
    =
    \min\left\{1,
    \exp\left[
    \log\pi(c'\mid z)-\log\pi(c\mid z)
    +\log q(c\mid c')
    \right]
    \right\}.
    \label{eq:merge-acceptance}
\end{equation}

\section{Acceptance-rate monitoring}

For each Metropolis--Hastings block, monitor empirical acceptance rates:
\begin{align}
    \widehat{\mathrm{AR}}_\theta
    &=\frac{\#\{\theta\text{ proposals accepted}\}}
            {\#\{\theta\text{ proposals attempted}\}},\\
    \widehat{\mathrm{AR}}_\beta
    &=\frac{\#\{\beta\text{ proposals accepted}\}}
            {\#\{\beta\text{ proposals attempted}\}},\\
    \widehat{\mathrm{AR}}_a
    &=\frac{\#\{a\text{ proposals accepted}\}}
            {\#\{a\text{ proposals attempted}\}},\\
    \widehat{\mathrm{AR}}_b
    &=\frac{\#\{b\text{ proposals accepted}\}}
            {\#\{b\text{ proposals attempted}\}},\\
    \widehat{\mathrm{AR}}_{\mathrm{split}}
    &=\frac{\#\{\text{split proposals accepted}\}}
            {\#\{\text{split proposals attempted}\}},\\
    \widehat{\mathrm{AR}}_{\mathrm{merge}}
    &=\frac{\#\{\text{merge proposals accepted}\}}
            {\#\{\text{merge proposals attempted}\}}.
\end{align}
The single-site label update in \eqref{eq:c-update} is a Gibbs update and
therefore has acceptance probability one.  Very low split acceptance can imply
that the current item positions support only coarse clusters, that $S_0$ is too
loose/tight for the scale of $z_q$, or that the restricted Gibbs split proposal
is not sufficiently close to high-posterior split configurations.

\section{Posterior cluster summaries}

Because labels are not identifiable, cluster summaries should be
label-switching invariant.  The primary summary is the posterior similarity
matrix
\begin{equation}
    C_{qr}
    =
    P(c_q=c_r\mid Y),
    \qquad q,r=1,\dots,P.
\end{equation}
During MCMC, update online by
\begin{equation}
    C_{qr}^{\mathrm{raw}}
    \leftarrow
    C_{qr}^{\mathrm{raw}}+\ind(c_q=c_r)
\end{equation}
after burn-in.  At the end, divide by the number of saved iterations.

Also monitor
\begin{equation}
    K_+^{(t)}=|\{l:n_l^{(t)}>0\}|,
\end{equation}
the layer-scale diagnostic in \eqref{eq:layer-scale-diagnostic}, and the
split/merge acceptance rates.


\section{One full MCMC sweep}

\begin{algorithm}[H]
\caption{One MCMC sweep: LSIRM update and collapsed label update}
\small
\begin{algorithmic}[1]
\State \textbf{Input:} data $Y$, layer maps $q\mapsto(\ell(q),j(q))$, hyperparameters, proposal scales, $K^\star$, and $M_{\mathrm{SM}}$.
\State Update layer-specific nuisance parameters $\xi_\ell$, if present.
\For{$\ell=1,\dots,L$}
    \For{$i=1,\dots,n_\ell$}
        \State Propose $\theta_i^{(\ell)\prime}$ and accept/reject using \eqref{eq:theta-mh-ratio}.
        \State Propose $a_i^{(\ell)\prime}$ and accept/reject using \eqref{eq:a-mh-ratio}.
    \EndFor
    \For{$j=1,\dots,P_\ell$}
        \State Propose $\beta_j^{(\ell)\prime}$ and accept/reject using \eqref{eq:beta-mh-ratio}.
        \State Let $q=q(\ell,j)$ and propose $b_j^{(\ell)\prime}$.
        \State Accept/reject $b_j^{(\ell)\prime}$ using \eqref{eq:b-mh-ratio}, including the mixture prior term.
    \EndFor
    \State Apply deterministic centering or identification constraints if used.
    \State Update $\sigma_{a,\ell}^2$, $\sigma_{\theta,\ell}^2$, and $\sigma_{\beta,\ell}^2$ using \eqref{eq:sigma-a-update}--\eqref{eq:sigma-beta-update}, if sampled.
\EndFor
\State Construct global item positions $z_q=b_{j(q)}^{(\ell(q))}$ for $q=1,\dots,P$.
\State Compute and store layer-scale diagnostics $s_\ell$ and $R_{\mathrm{scale}}$.
\For{$q=1,\dots,P$}
    \State Remove item $q$ from its current cluster.
    \State Compute NIW leave-one-out predictive weights in \eqref{eq:c-update} for $l=1,\dots,K^\star$.
    \State Sample $c_q$ from the normalized weights.
\EndFor
\end{algorithmic}
\end{algorithm}

\begin{algorithm}[H]
\caption{One MCMC sweep: split--merge, mixture-parameter update, and storage}
\small
\begin{algorithmic}[1]
\For{$m=1,\dots,M_{\mathrm{SM}}$}
    \State Select two anchors $u,v$ uniformly at random.
    \If{$c_u=c_v$}
        \State Attempt a restricted-Gibbs split move and accept/reject using \eqref{eq:split-acceptance}.
    \Else
        \State Attempt a merge move and accept/reject using \eqref{eq:merge-acceptance}.
    \EndIf
\EndFor
\For{$l=1,\dots,K^\star$}
    \State Compute NIW posterior parameters $(\kappa_l,\nu_l,m_l,S_l)$ from \eqref{eq:niw-posterior-params}.
    \State Draw $\Sigma_l\mid z,c$ using \eqref{eq:sigma-posterior}.
    \State Draw $\mu_l\mid \Sigma_l,z,c$ using \eqref{eq:mu-posterior}.
\EndFor
\State Optionally draw $\rho\mid c$ using \eqref{eq:rho-posterior}.
\If{iteration is after burn-in}
    \State Store $c$, $K_+$, $z$, selected LSIRM parameters, acceptance indicators, and layer-scale diagnostics.
    \State Update posterior similarity counts $C_{qr}^{\mathrm{raw}}\leftarrow C_{qr}^{\mathrm{raw}}+\ind(c_q=c_r)$.
\EndIf
\end{algorithmic}
\end{algorithm}

\end{document}
